%% kapa: exam=1 round=37 year=2026 period=2 form=A subject=appraisal_law subject_ko=감정평가관계법규 seq=34 number=34
%% source: Q-Net 감정평가사 기출문제 (cst00309) artlSeq=5258924 file=raw/1008/5258924/2026년 제37회 감정평가사 1차 2교시.pdf
%% conversion: kordoc 4.24.51 -> tools/gen1.py (auto); review=proofread-sonnet; flags=none
\begin{question}{34}
부동산등기법령상 이의에 관한 설명으로 옳지 않은 것은?
\begin{choices}
  \item 새로운 사실이나 새로운 증거방법을 근거로 이의신청을 할 수 있다.
  \item 등기관은 이의가 이유 없다고 인정하면 이의신청일부터 3일 이내에 의견을 붙여 이의신청서 또는 이의신청정보를 관할 지방법원에 보내야 한다.
  \item 이의에는 집행정지의 효력이 없다.
  \item 관할 지방법원은 이의신청에 대하여 결정하기 전에 등기관에게 가등기 또는 이의가 있다는 뜻의 부기등기를 명령할 수 있다.
  \item 송달에 대하여는 ｢민사소송법｣을 준용하고, 이의의 비용에 대하여는 ｢비송사건절차법｣을 준용한다.
\end{choices}
\end{question}
\answer{1}
